\chapter{Forms and Characteristic Boundaries}
\label{ch:real-forms}

The complex classification does not identify a real algebra merely from
its Dynkin diagram. A real form also specifies a conjugation. We first
translate conjugations into involutions of a compact form, then solve
the resulting finite classification problem. The final section gives
counterexamples to some hypotheses used earlier.

\section{Automorphisms and descent}

\begin{theorem}[Automorphisms of a complex semisimple algebra]\label{res:J01}
For a complex semisimple algebra with Dynkin diagram \(D\), including
its connected components,
\[
 \operatorname{Aut}(\mathfrak g)
 =\operatorname{Int}(\mathfrak g)\rtimes\operatorname{Aut}(D).
\]
Here \(\operatorname{Int}\) is generated by exponentials of inner
derivations; automorphisms of \(D\) may permute isomorphic components.
\end{theorem}
\begin{proof}
A permutation of the simple generators preserving the Cartan matrix
preserves every Serre relation, so gives a Lie automorphism by
Theorem~\ref{res:F03}. It gives a section of the asserted quotient.
For an arbitrary automorphism, Cartan conjugacy first lets us compose
with an inner automorphism to preserve a fixed Cartan. A Weyl-group
representative then makes it preserve the positive roots. Its action
on simple roots is an automorphism of \(D\).
Remove that diagram automorphism.

The remaining map fixes the Cartan pointwise, sends each simple
root vector \(e_i\) to \(c_i e_i\), and sends \(f_i\) to
\(c_i^{-1}f_i\), because \([e_i,f_i]=h_i\).
Choose \(h\) in the Cartan with \(\exp(\alpha_i(h))=c_i\):
choose complex logarithms of the \(c_i\), and use that the
simple roots are a basis. Then \(\exp(\operatorname{ad}h)\)
has precisely these effects. The simple generators generate the
algebra, so the automorphism is this inner one.

Finally a nontrivial diagram automorphism permutes at least two
fundamental weights (also when it permutes simple factors).
The corresponding finite-dimensional simple modules are not
isomorphic by Theorem~\ref{res:G06}. Twisting by an inner
automorphism, however, preserves every module's isomorphism class:
for \(\exp(\operatorname{ad}x)\), the intertwiner is
\(\exp(\rho(x))\), as the exponential power series gives
\(\rho(\exp(\operatorname{ad}x)y)
=\exp(\rho(x))\rho(y)\exp(-\rho(x))\);
compose these intertwiners for a product of inner exponentials.
Thus no nontrivial diagram automorphism is inner. The section
has trivial intersection with the inner subgroup, proving the
semidirect product.
See \cite[Section 14.3]{humphreys} and
\cite[Proposition 40.3]{etingof}.
\end{proof}

\begin{proposition}[Finite Galois descent and characteristic-zero Ado]\label{res:J02}
Every finite reduced root system has a split semisimple model over
every characteristic-zero field, obtained from its Chevalley
integer lattice.
Let \(L/K\) be a finite Galois extension and let \(\mathfrak a\) be
a Lie algebra over \(K\). The \(K\)-forms of
\(\mathfrak a_L=L\otimes_K\mathfrak a\) are classified by
semilinear Galois descent data, or equivalently by
\[
 H^1(\operatorname{Gal}(L/K),\operatorname{Aut}_L(\mathfrak a_L)).
\]
The equivalence is by coboundary, not by equality of cocycles.
In addition, every finite-dimensional Lie algebra over an arbitrary
characteristic-zero field has a faithful finite-dimensional representation.
\end{proposition}
\begin{proof}
Start with the Chevalley basis of Theorem~\ref{res:F05}. Its
integer span \(\mathfrak g_{\mathbb Z}\) is closed under brackets.
All coefficients of its Jacobi and alternation identities are
integers and are zero after embedding in \(\mathbb C\), hence are
zero already over \(\mathbb Z\). Thus
\(\mathfrak g_K=K\otimes_{\mathbb Z}\mathfrak g_{\mathbb Z}\)
is a Lie algebra with that same basis over any field \(K\).
Its Killing matrix has integer entries, because each adjoint
matrix does. Its determinant is a nonzero integer: over
\(\mathbb C\) it is the nondegenerate Killing form of the
constructed semisimple algebra. In characteristic zero this
determinant remains nonzero, so \(\mathfrak g_K\) is semisimple
by Theorem~\ref{res:B05}. Its Cartan basis \(h_i\) acts diagonally
with the integer eigenvalues \(\alpha(h_i)\); the root spaces
are already the individual \(Kx_\alpha\), and the zero space
is \(Kh_1+\cdots+Kh_r\). This is the asserted split model.
The descent below applies in particular to its forms trivialized
over the stated extension \(L\); it does not identify all
nonsplit forms with the split model itself.

Write \(\Gamma=\operatorname{Gal}(L/K)\), and let \(\gamma_0\)
act on scalar coefficients. A semilinear bracket-preserving action
has the form \(\widetilde\gamma=c_\gamma\gamma_0\).
The action law is exactly
\(c_{\gamma\delta}=c_\gamma\gamma_0(c_\delta)\).
Changing an identification with \(\mathfrak a_L\) by \(b\)
replaces it by
\(c'_\gamma=b c_\gamma\gamma_0(b)^{-1}\).

We recall why the fixed subspace actually descends the vector space.
Distinct automorphisms of a field are linearly independent as functions:
in a minimal relation, compare the relation at \(au\) with that at
\(u\) multiplied by one automorphism's value at \(a\); choosing \(a\)
on which two automorphisms differ gives a shorter relation.
For a nonzero semilinear \(\Gamma\)-space \(V\), the expression
\(\sum_{\gamma\in\Gamma}\widetilde\gamma(a v)\) is invariant.
For some \(a,v\) it is nonzero, by that independence, applied
after an \(L\)-linear coordinate functional.
More generally, let \(W\) be the \(L\)-span of all invariant
vectors. If \(V/W\ne0\), the same expression computed in the
quotient produces a nonzero invariant sum that lifts to an
invariant vector of \(V\), a contradiction. Thus \(W=V\).
Choose an \(L\)-basis from the invariant vectors. An invariant
linear combination in that basis has all its coefficients in
\(K\). Hence
\(L\otimes_K V^\Gamma\to V\) is an isomorphism.
Taking fixed points of the bracket gives the asserted descended
Lie algebra, and recovers all forms. Any finite-dimensional form
that splits over an algebraic closure already splits over some
finite Galois extension: finitely many entries of an isomorphism
and its inverse generate a finite extension, which can be enlarged
to its Galois closure in characteristic zero.

Here is the additional descent needed for Ado; no embedding of the
whole original field into \(\mathbb C\) is required.
Let \(F\subseteq K\) be the finitely generated extension of
\(\mathbb Q\) generated by the finitely many structure constants
in a chosen basis of \(\mathfrak g\). The field \(F\) embeds into
\(\mathbb C\): send a transcendence basis to algebraically
independent complex numbers and extend through its finite
algebraic extension. By Theorem~\ref{res:I11}, the corresponding
complex algebra has a faithful representation of some dimension
\(n\). Unknown matrices \(X_i\in M_n\) for the basis elements
satisfy finitely many equations
\[
 [X_i,X_j]=\sum_\ell c_{ij}^{\ell}X_\ell.
\]
Faithfulness is the nonvanishing of one maximal minor \(\Delta\)
of the matrix whose columns are the entries of the \(X_i\).
Adjoin a variable \(u\) and the equation \(u\Delta-1=0\).
This polynomial system over \(F\) has a complex solution and
therefore its ideal is proper. A maximal ideal containing it
has residue field finite algebraic over \(F\), by Zariski's lemma.
For completeness, that lemma follows by choosing a transcendence
basis in a field that is finitely generated as an algebra:
clearing the denominators of the remaining algebraic generators
would make the field integral over
\(F[t_1,\ldots,t_r,1/f]\). An integral subring of a field over
which the field is integral must itself be a field, whereas for
\(r>0\) the displayed ring misses the inverse of some \(t_1-a\),
\(a\in F\). Thus \(r=0\), and the algebraic extension is finite.

Embed this residue field in an algebraic closure of \(K\).
Its compositum with \(K\) is a finite extension \(K'/K\), and
the resulting matrices are a faithful representation of
\(\mathfrak g_{K'}\). Restrict the scalars of the representation
space to \(K\). A nonzero original vector remains nonzero after
base extension, so the \(K\)-representation is faithful.
This argument explains why merely collecting finitely many
complex coefficients, without the polynomial-system step,
would not justify descent.
For forms compare \cite[Proposition 40.4]{etingof}.
\end{proof}

\section{The compact form and its involutions}

\begin{theorem}[Conjugations and the compact form]\label{res:J03}
A real form of a complex Lie algebra is the fixed algebra of an
antilinear bracket-preserving involution, and its complexification
is the original algebra. Every complex semisimple algebra has
a compact real form \(\mathfrak u\), on which the Killing form
is negative definite.
\end{theorem}
\begin{proof}
For an antilinear involution \(\sigma\), write
\[
 x=\frac{x+\sigma x}{2}
   +i\,\frac{x-\sigma x}{2i}.
\]
Both displayed summands before multiplication by \(i\) are fixed.
Their real direct sum is the whole complex vector space,
so complexification of the fixed algebra is an isomorphism.
Conversely the usual coefficient conjugation in a complexification
has exactly the original real algebra as fixed points.
Isomorphisms of real forms are precisely complex automorphisms
intertwining their conjugations.

Take the Chevalley normalization of Theorem~\ref{res:F05}
(equivalently its compact-normalization construction):
\([e_\alpha,e_{-\alpha}]=h_\alpha\),
\(B(e_\alpha,e_{-\alpha})>0\), and the conjugate-linear rule
\[
 \omega(h_i)=-h_i,\qquad
 \omega(e_i)=-f_i,\qquad \omega(f_i)=-e_i
\]
preserves the defining Serre relations. It is an involution.
Choose the root vectors compatibly so
\(\omega(e_\alpha)=-e_{-\alpha}\).
The fixed real space has basis
\[
 ih_i,\qquad
 e_\alpha-e_{-\alpha},\quad i(e_\alpha+e_{-\alpha})
 \quad(\alpha>0).
\]
It is a Lie algebra because it is a fixed algebra.
On the real span of the coroots the Killing form is positive
definite, by the positive-definite root form. Multiplication by
\(i\) makes it negative. Different pairs
\(\{\alpha,-\alpha\}\) are orthogonal, and on each of the two
displayed real root vectors its value is
\(-2B(e_\alpha,e_{-\alpha})<0\); their cross term is zero.
Thus \(B|_{\mathfrak u}\) is negative definite.
This construction uses the root normalization proved earlier,
not the existence of a compact group to prove that normalization.
Compare \cite[Proposition 41.1]{etingof} and
\cite[Section VI.1]{knapp}.
\end{proof}

\begin{theorem}[Polar reduction of real structures]\label{res:J04}
Fix the compact conjugation \(\omega\). Real forms are in bijection
with conjugacy classes of involutions
\(\theta\in\operatorname{Aut}_{\mathbb R}(\mathfrak u)\), by
\(\theta\mapsto\sigma=\theta\omega\).
If \(\mathfrak u=\mathfrak u^+\oplus\mathfrak u^-\) are its
eigenspaces, the corresponding real form is
\[
 \mathfrak g_\theta=\mathfrak u^+\oplus i\mathfrak u^-.
\]
The fixed complex algebra \(\mathfrak k=\mathfrak g^\theta\)
is reductive. The Killing signature of the real form, recorded
as positive dimension minus negative dimension, is
\(\dim_{\mathbb C}\mathfrak g-2\dim_{\mathbb C}\mathfrak k\).
\end{theorem}
\begin{proof}
Use the positive Hermitian form
\(\langle x,y\rangle=-B(x,\omega y)\).
Invariance of \(B\) gives, for every complex automorphism \(a\),
\[
 a^\dagger=\omega a^{-1}\omega,\qquad
 \omega(a)=(a^\dagger)^{-1}.
\]
Write any real structure as \(\sigma=a\omega\).
The condition \(\sigma^2=1\) says \(a\omega(a)=1\), hence
\(a=a^\dagger\). Its nonzero eigenvalues are real, and
\([V_\beta,V_\gamma]\subseteq V_{\beta\gamma}\).
Therefore \(|a|^t\), which multiplies \(V_\beta\) by
\(|\beta|^t\), is a Lie automorphism for every real \(t\):
the scalar on a bracket is the product of its scalars on the
two entries. Put \(\theta=a|a|^{-1}\).
It has eigenvalues \(\pm1\), is unitary, and commutes with
\(\omega\). Conjugating \(a\omega\) by \(|a|^{-1/2}\) gives
\(\theta\omega\). This proves existence of the reduction
without assuming that an arbitrary matrix power preserves brackets.

Suppose \(\theta\omega\) and \(\theta'\omega\) are intertwined
by \(x\). Thus \(\theta'=x\theta\omega(x)^{-1}=x\theta x^\dagger\).
Writing \(P=x^\dagger x\), the equation \((\theta')^2=1\)
gives \(\theta P\theta=P^{-1}\).
The operator \(v=xP^{-1/2}\) is a unitary Lie automorphism,
so it preserves \(\mathfrak u\), and
\[
 v\theta v^{-1}=xP^{-1/2}\theta P^{1/2}x^{-1}
              =x\theta P x^{-1}=\theta'.
\]
The converse is immediate. This proves the asserted bijection.

The real fixed algebra \(\mathfrak u^+\) has the invariant positive
form \(-B\). In any compact algebra with such a form, the
orthogonal complement of an ideal is an ideal. Moreover an
abelian ideal is central: for \(a,b\) in it and arbitrary \(x\),
\(\langle[x,a],b\rangle=\langle x,[a,b]\rangle=0\), while
\([x,a]\) still belongs to the ideal. Taking a last nonzero
derived ideal of the radical, splitting its orthogonal complement,
and inducting shows that the entire radical is central.
Thus the algebra, and its complexification \(\mathfrak k\),
are reductive by Theorem~\ref{res:I07}.
Finally the two eigenspaces are \(B\)-orthogonal;
\(B\) is negative on \(\mathfrak u^+\) and positive on
\(i\mathfrak u^-\). This gives the signature formula.
See \cite[Theorem 41.5]{etingof}.
\end{proof}

\begin{proposition}[A stable Cartan and Vogan data]\label{res:J05}
An involution of a compact form has a stable complex Cartan
\(\mathfrak h=\mathfrak h^+\oplus\mathfrak h^-\), with
\(\mathfrak h^+\) a Cartan in \(\mathfrak k=\mathfrak g^\theta\),
and no root vanishes identically on \(\mathfrak h^+\).
One can choose positive roots preserved by \(\theta\).
Thus \(\theta\) is encoded by an involution of the Dynkin
diagram, with a sign \(+1\) or \(-1\) on each fixed vertex.
Paired vertices need no additional sign.
\end{proposition}
\begin{proof}
Choose a maximal abelian subalgebra \(\mathfrak t^+\) of
\(\mathfrak u^+\). Choose a maximal subspace
\(\mathfrak t^-\subseteq\mathfrak u^-\) that commutes with
\(\mathfrak t^+\) and with itself. Then
\(\mathfrak t=\mathfrak t^+\oplus\mathfrak t^-\) is maximal
abelian in \(\mathfrak u\). Indeed an element centralizing it
has a positive and a negative part under \(\theta\); maximality
puts those parts in \(\mathfrak t^+\) and \(\mathfrak t^-\).
All its adjoint operators are commuting skew-adjoint operators,
so \(\mathfrak h=\mathfrak t_{\mathbb C}\) is toral and
self-centralizing, hence Cartan. The same argument in
\(\mathfrak u^+\) proves that
\(\mathfrak h^+=\mathfrak t^+_{\mathbb C}\) is Cartan in
\(\mathfrak k\).

If a root \(\alpha\) vanished on \(\mathfrak h^+\), its
coroot would lie in \(\mathfrak h^-\), using the
\(\theta\)-invariant root inner product. Hence
\(\theta\alpha=-\alpha\). For nonzero root vectors,
\(\theta(e_\alpha)=c e_{-\alpha}\) and
\(\theta(e_{-\alpha})=c^{-1}e_\alpha\).
The vector \(e_\alpha+c e_{-\alpha}\) would be fixed by
\(\theta\), would centralize \(\mathfrak h^+\), and would
not belong to it. This contradicts its being Cartan in
the reductive \(\mathfrak k\).

Choose a generic element in \(i\mathfrak t^+\), avoiding
the finitely many root hyperplanes. It orders the roots and
is fixed by \(\theta\), so the resulting positive system is
preserved. On simple generators
\(\theta h_i=h_{\sigma(i)}\).
If \(\sigma(i)=i\), involutivity forces
\(\theta e_i=\varepsilon_i e_i\),
\(\theta f_i=\varepsilon_i f_i\), \(\varepsilon_i=\pm1\).
If \(i,\sigma(i)\) are paired, choose \(e_{\sigma(i)}=\theta e_i\)
and the compatible negative root vector; the second application
of \(\theta\) is then the identity. Conversely every such
assignment respects the Serre relations, commutes with compact
conjugation, and squares to one.
These signs and pairs are the Vogan data.
See \cite[Proposition 41.7 and Lemma 41.8]{etingof}.
\end{proof}

\section{Equivalence and the classical list}

\begin{proposition}[Equivalence of the finite data]\label{res:J06}
Two normalized Vogan data define isomorphic real forms exactly when
their involutions are conjugate by a compact Cartan normalizer,
after possibly moving one stable Cartan by the connected fixed group.
Thus the test is finite: test root-system automorphisms and normalize
the scalar multiples on paired root vectors. In the inner case it
is the action of \(W\rtimes\operatorname{Aut}(D)\) on
\(\operatorname{Hom}(Q,\{\pm1\})\), where \(Q\) is the root lattice.
Writing a sign as \((-1)^{b_j}\), the simple reflection at \(i\) acts by
\[
 b_j\longmapsto b_j-a_{ij}b_i\pmod2,\qquad
 a_{ij}=\alpha_j(h_i).
 \tag{10.1}\label{eq:vogan-reflection}
\]
In particular a black fixed vertex changes the colors of its
fixed neighbors for which \(a_{ij}\) is odd.
\end{proposition}
\begin{proof}
We first justify the Cartan adjustment in the assertion.
Maximal abelian subalgebras of a compact reductive algebra
\(\mathfrak v\) are conjugate under its connected inner group.
Here is a useful finite-dimensional proof. Let \(H\) be the
identity component of the group of orthogonal automorphisms
of \(\mathfrak v\) that fix \(Z(\mathfrak v)\) pointwise.
It is a closed subgroup of an orthogonal group and hence compact.
Its Lie algebra consists of skew-adjoint derivations that vanish
on the center. Theorem~\ref{res:B07}, applied to the semisimple
summand of \(\mathfrak v_{\mathbb C}\), identifies these
derivations with \(\operatorname{ad}\mathfrak v\).
Choose regular elements \(X,Y\)
in two maximal abelian subalgebras; regular means that their
centralizers are those subalgebras, obtained by avoiding finitely
many weight hyperplanes. Maximize
\(\langle\operatorname{Ad}(g)X,Y\rangle\) in the compact group.
Differentiation in the direction \(Z\) gives
\(\langle Z,[\operatorname{Ad}(g)X,Y]\rangle=0\) for every \(Z\).
Thus \(\operatorname{Ad}(g)X\) commutes with \(Y\), and belongs
to the second maximal abelian subalgebra. Its centralizer is
the first subalgebra transported by \(g\), so the two maximal
abelian subalgebras coincide after transport.
Every element of \(H\) is a finite product of inner
exponentials. This follows by the inverse
function theorem applied to a product of exponentials of a
basis at the identity; their image contains a neighborhood,
and a subgroup containing a neighborhood is open and equals
the identity component.

Suppose \(c\theta c^{-1}=\theta'\) in
\(\operatorname{Aut}(\mathfrak u)\).
It carries \(\mathfrak t^+\) to a maximal abelian subalgebra
of \(\mathfrak u^{\theta'}\).
Apply the preceding result in this fixed compact algebra to
make \(c\mathfrak t^+=(\mathfrak t')^+\).
The required correcting automorphisms are exponentials of
elements of \(\mathfrak u^{\theta'}\), so commute with \(\theta'\).
By Proposition~\ref{res:J05}, no root vanishes on these
positive Cartan parts. Their centralizers in \(\mathfrak g\)
are consequently the full stable Cartans. Thus \(c\) can
be made to carry the full Cartan to the full Cartan.

The normalizer acts on roots through a finite root-system
automorphism. Its kernel scales each simple root vector
by a unit scalar and is the compact torus: solve the
simple-root logarithms as in Theorem~\ref{res:J01}.
There is therefore an explicit finite equivalence test.
For each root-system automorphism \(w\), use its lift formed
from root \(\mathfrak{sl}_2\) reflections and diagram permutations.
Retain it only if it conjugates the diagram action of the
first datum to that of the second. Compare the resulting
multipliers on the simple root vectors. Conjugation by the
kernel torus changes a multiplier by \(z_j/z_{\sigma(j)}\).
It cannot change a fixed-vertex sign, and it can normalize
both multipliers on any paired vertices, since they are
reciprocals. Thus comparison of the fixed signs is sufficient
after this finite set of tests.

For an inner datum the involution fixes the Cartan and is
the character \(b\) on \(Q\). A reflection changes its value
on \(\alpha_j\) to its value on
\(\alpha_j-a_{ij}\alpha_i\), proving
\eqref{eq:vogan-reflection}. It does not change \(b_i\), since
\(a_{ii}=2\). The same calculation at a fixed simple root in
an outer diagram gives the stated fixed-neighbor move,
after renormalizing paired generators.
This proves both the conjugacy criterion and the elementary
moves used below. Compare \cite[Lemma 42.1]{etingof}.
\end{proof}

\begin{theorem}[Classical simple real Lie algebras]\label{res:J07}
Up to the low-rank identifications specified below, the real forms
of the simple classical complex algebras are the following.
Parameters \(p,q\) satisfy \(p+q=n\), except in the orthogonal
row where \(p+q=N\); interchange of \(p,q\) gives the same form.
\[
\begin{array}{c|l|l}
 \mathfrak g&\text{real forms}&
       (\text{maximal compact algebra})_{\mathbb C}\\ \hline
 \mathfrak{sl}_n&
 \mathfrak{su}(p,q)&
 \mathfrak s(\mathfrak{gl}_p\oplus\mathfrak{gl}_q)\\
 &\mathfrak{sl}_n(\mathbb R)&\mathfrak{so}_n\\
 &\mathfrak{su}^*(2m),\ n=2m&
             \mathfrak{sp}_{2m}\\ \hline
 \mathfrak{so}_N&
 \mathfrak{so}(p,q)&\mathfrak{so}_p\oplus\mathfrak{so}_q\\
 &\mathfrak{so}^*(2n),\ N=2n&\mathfrak{gl}_n\\ \hline
 \mathfrak{sp}_{2n}&
 \mathfrak{sp}(p,q)&\mathfrak{sp}_{2p}\oplus\mathfrak{sp}_{2q}\\
 &\mathfrak{sp}_{2n}(\mathbb R)&\mathfrak{gl}_n
\end{array}
\]
The compact cases are included by \(q=0\).
The quaternionic row \(\mathfrak{su}^*(2m)\) is additional for
\(m\geq2\).
Every complex simple algebra, regarded as a real algebra, is
also real simple; its complexification has two isomorphic
simple components. These form the other branch of the
classification of real simple algebras.
\end{theorem}
\begin{proof}
The proof classifies involutions of the compact forms from
Theorem~\ref{res:J04}; all reductions used here are ordinary
spectral or bilinear-form calculations.

For type \(A_{n-1}\), an inner involution of
\(\mathfrak{su}(n)\) is \(\operatorname{Ad}A\).
Commutation with compact conjugation says \(A^\dagger A\)
is scalar, so scale \(A\) to be unitary.
Its square is scalar; multiplying by a unit scalar makes
\(A^2=1\). The Hermitian spectral theorem puts it in the
form \(I_p\oplus(-I_q)\).
The real structure on \(\mathfrak{sl}_n\) is
\(X\mapsto-A X^\dagger A\); its fixed algebra is exactly
\(\mathfrak{su}(p,q)\).
Its compact fixed algebra consists of the block-diagonal
matrices, giving the first row of the table.

For \(n>2\) the nontrivial diagram action is represented by
\(X\mapsto-X^T\). An outer involution therefore has the form
\(-A X^T A^{-1}\). Its square is inner by \(A A^{-T}\);
this is scalar exactly when \(A^T=\varepsilon A\), with
\(\varepsilon=\pm1\). In the compact case \(A\) can be
chosen unitary. A unitary symmetric matrix is congruent
to \(I\): its associated antiunitary map squares to \(1\),
whose fixed real orthonormal basis is obtained by writing
each vector as fixed real and imaginary parts and applying
real Gram--Schmidt. A unitary skew-symmetric matrix is
congruent to
\(J=\left(\begin{smallmatrix}0&I_m\\-I_m&0\end{smallmatrix}\right)\):
for its antiunitary map \(j\), choose a unit vector \(v\);
then \(v,jv\) are orthogonal because \(j^2=-1\), split off
this pair, and induct. Thus the skew case exists only
when \(n=2m\).
The resulting conjugations are \(X\mapsto\overline X\)
and \(X\mapsto J\overline XJ^{-1}\).
Their fixed algebras are
\(\mathfrak{sl}_n(\mathbb R)\) and
\(\mathfrak{su}^*(2m)=\mathfrak{sl}_m(\mathbb H)\), where
the latter notation means quaternionic matrices with real
part of the trace zero. The compact fixed subalgebras
of the involutions are \(\mathfrak{so}_n\) and
\(\mathfrak{sp}_{2m}\), by their defining matrix equations.

For orthogonal type, every compact involution is conjugate by a
compact automorphism to an involution represented by
\(A\in O(N,\mathbb R)\). By Theorem~\ref{res:J01}, the complex
inner subgroup is represented by \(SO(N,\mathbb C)\), and
in even dimension a real determinant \(-1\) representative
induces a diagram transposition. In \(D_4\) the diagram action
of an involution is either trivial or conjugate to this
transposition in the triality group \(S_3\). The permutations
of compact-normalized simple generators lift the required
diagram conjugacies, since they commute with compact conjugation
on those generators. Thus the compact involution, after this
conjugacy if needed, has a complex orthogonal representative \(A\).
As it preserves the compact real algebra, \(\overline A=cA\)
for a scalar \(c\), and orthogonality gives \(c=\pm1\).
The case \(c=-1\) would make \(A=iR\) with \(R\) real and
\(R^TR=-I\), which is impossible. Therefore \(A\) is real.
The condition \((\operatorname{Ad}A)^2=1\) is \(A^2=\pm I\).
If \(A^2=I\), the real orthogonal eigenspaces give
\(A=I_p\oplus(-I_q)\). The fixed real vector space of
the antiunitary \(A\) times coefficient conjugation is
\(\mathbb R^p\oplus i\mathbb R^q\); the original symmetric
form restricts to signature \((p,q)\). Hence the fixed
Lie algebra is \(\mathfrak{so}(p,q)\).
If \(A^2=-I\), \(N=2n\) and a real orthonormal basis
puts \(A\) into the standard complex-structure matrix
\(J\). This gives \(\mathfrak{so}^*(2n)\), the algebra
preserving both the complex symmetric form and the
quaternionic structure \(J\) times conjugation. Its
compact centralizer is \(\mathfrak u(n)\).
These are all orthogonal cases. A matrix with square
\(-I\) has determinant \(1\), so no missing
determinant-\(-1\) case occurs.

For symplectic type all diagram automorphisms are trivial.
An involution of the compact algebra is represented by
\(A\in Sp(n)\), modulo its central \(\{\pm I\}\);
thus \(A^2=\pm I\). For \(A^2=I\), the eigenspaces
are quaternionic orthogonal subspaces of dimensions
\(p,q\). They give
\(\mathfrak{sp}(p,q)\) and compact centralizer
\(\mathfrak{sp}(p)\oplus\mathfrak{sp}(q)\).
For \(A^2=-I\), a quaternionic orthonormal basis makes
\(A=iI\): choose a complex eigenvector of eigenvalue
\(i\), whose quaternionic mate has eigenvalue \(-i\),
and induct on its orthogonal complement.
The compact centralizer is \(U(n)\).
If \(j\) denotes the compact quaternionic structure on
\(\mathbb C^{2n}\), then the new antiunitary map \(Aj\)
squares to \(+I\). Its fixed real space has a real
nondegenerate alternating form, so a real symplectic
basis identifies the fixed algebra with
\(\mathfrak{sp}_{2n}(\mathbb R)\).
This also proves the last two entries of the table directly.

Within each family the matrix normal forms are unique up
to the stated exchange of eigenspaces. Distinct entries
are separated by the compact fixed algebra and its
center, except for the following actual root identifications.
One has \(B_1=C_1=A_1\), with compact
\(\mathfrak{so}(3)\cong\mathfrak{su}(2)\) and split
\(\mathfrak{so}(2,1)\cong\mathfrak{sl}_2(\mathbb R)\).
The isomorphism \(B_2=C_2\) identifies
\(\mathfrak{so}(5),\mathfrak{so}(4,1),\mathfrak{so}(3,2)\)
with \(\mathfrak{sp}(2),\mathfrak{sp}(1,1),
\mathfrak{sp}_4(\mathbb R)\).
For \(D_3=A_3\) the five forms correspond as
\[
\begin{array}{c|ccccc}
 &\mathfrak{so}(6)&\mathfrak{so}(5,1)&\mathfrak{so}(4,2)&
       \mathfrak{so}(3,3)&\mathfrak{so}^*(6)\\ \hline
 &\mathfrak{su}(4)&\mathfrak{su}^*(4)&\mathfrak{su}(2,2)&
       \mathfrak{sl}_4(\mathbb R)&\mathfrak{su}(3,1).
\end{array}
\]
These identifications follow by transporting their compact
involutions under the indicated root-system isomorphisms.
For \(D_4\) take the simple roots \(e_1-e_2,e_2-e_3,
e_3-e_4,e_3+e_4\) for rotations on four ordered planes.
The compact involution for \(\mathfrak{so}^*(8)\) has
angles \((\pi/2,\pi/2,\pi/2,\pi/2)\), so the eigenvalue
on the root line for the last simple root is \(-1\)
and that on the other three is \(1\). The involution for
\(\mathfrak{so}(6,2)\) has angles \((\pi,0,0,0)\), so
only the first simple root line has eigenvalue \(-1\).
Triality exchanges these two end roots. The eigenvalues
on the simple root lines determine those on all root
lines, and each automorphism acts trivially on the
common torus. This identifies the compact involutions
and therefore identifies their real forms.
The nonsimple case \(D_2=A_1+A_1\) has
\(\mathfrak{so}(4)=\mathfrak{su}(2)\oplus\mathfrak{su}(2)\),
\(\mathfrak{so}(2,2)=\mathfrak{sl}_2(\mathbb R)^{\oplus2}\),
and \(\mathfrak{so}(3,1)=\mathfrak{sl}_2(\mathbb C)\)
as a real algebra. The remaining mixed form is
\(\mathfrak{so}^*(4)\cong
\mathfrak{su}(2)\oplus\mathfrak{sl}_2(\mathbb R)\):
its compact fixed algebra is
\(\mathfrak{su}(2)\oplus\mathfrak u(1)\), so the two
\(A_1\) factors have respectively compact and split involutions.
The usual simple ranks are
\(A_{n-1}\) for \(n\geq2\), \(B_n,C_n\) for \(n\geq2\),
and \(D_n\) for \(n\geq4\), with these identifications understood.

Finally let \(\mathfrak l\) be real simple. Its complexification
is semisimple: a nonzero radical would be conjugation-stable
and would descend to a solvable ideal. Conjugation permutes
its complex simple ideals transitively; otherwise a union
of orbits descends to a proper real ideal. An involution
has orbits of size one or two, so either the complexification
is simple, or it is a pair \(\mathfrak a\oplus\mathfrak a\)
with factors exchanged. In the latter case the fixed algebra
is \(\mathfrak a\) with scalars restricted to \(\mathbb R\).
Conversely an ideal in such a real algebra complexifies
to a conjugation-stable ideal in the pair and hence is
zero or the whole algebra. This proves the second branch.
The matrix classification agrees with
\cite[Section 41.4]{etingof} and
\cite[Theorem 6.105]{knapp}.
\end{proof}

\section{Exceptional real forms by a finite root calculation}

The following conventions specify the finite calculation completely.
For \(E_r\), \(r=6,7,8\), label the horizontal chain
\(1-2-3-\cdots-(r-1)\), and attach the last vertex \(r\)
to vertex \(3\). Thus the three arms have lengths
\((2,r-4,1)\). For \(F_4\), vertices \(1,2\) are short,
vertices \(3,4\) long, and
\[
 A_{F_4}=\begin{pmatrix}
 2&-1&0&0\\-1&2&-2&0\\0&-1&2&-1\\0&0&-1&2
 \end{pmatrix},\qquad
 A_{G_2}=\begin{pmatrix}2&-3\\-1&2\end{pmatrix}.
\]
All \(E_r\) off-diagonal entries on an edge are \(-1\).
The convention remains \(a_{ij}=\alpha_j(h_i)\).

\begin{theorem}[The exceptional list]\label{res:J08}
The inner involutions, and the complexifications of their compact
fixed algebras, are exactly those in the following table.
\(S\) is the set of black vertices in the specified numbering,
and the orbit size refers to all \(2^r\) colorings.
The signature \(s\) is the difference of the two Killing indices.
\[
\begin{array}{c|c|r|l|r|r}
 D&S&\text{orbit}&\mathfrak k&\dim\mathfrak k&s\\ \hline
 G_2&\varnothing&1&G_2&14&-14\\
 G_2&\{1\}&3&A_1+A_1&6&2\\ \hline
 F_4&\varnothing&1&F_4&52&-52\\
 F_4&\{1\}&3&B_4&36&-20\\
 F_4&\{3\}&12&C_3+A_1&24&4\\ \hline
 E_6&\varnothing&1&E_6&78&-78\\
 E_6&\{1\}&27&D_5+\mathbb C&46&-14\\
 E_6&\{2\}&36&A_5+A_1&38&2\\ \hline
 E_7&\varnothing&1&E_7&133&-133\\
 E_7&\{1\}&63&D_6+A_1&69&-5\\
 E_7&\{4\}&36&A_7&63&7\\
 E_7&\{2,4\}&28&E_6+\mathbb C&79&-25\\ \hline
 E_8&\varnothing&1&E_8&248&-248\\
 E_8&\{1\}&135&D_8&120&8\\
 E_8&\{2\}&120&E_7+A_1&136&-24
\end{array}
\]
For \(E_6\) there are additionally two outer classes. The all-white
diagram involution has \(\mathfrak k=F_4\), dimension \(52\),
and signature \(-26\). Its nonzero fixed-vertex color class
has \(\mathfrak k=C_4\), dimension \(36\), and signature \(6\).
No other exceptional real forms occur.
The usual notation \(E_{r(s)}\), \(F_{4(s)}\), or \(G_{2(s)}\)
uses the signature column as its subscript.
\end{theorem}
\begin{proof}
For an inner datum \(b\in\mathbb F_2^r\), its fixed algebra is
\[
 \mathfrak h\oplus
 \bigoplus_{\substack{\alpha=\sum c_i\alpha_i\\
                     \sum b_i c_i\equiv0\ (2)}}\mathfrak g_\alpha.
 \tag{10.2}\label{eq:even-fixed-roots}
\]
Thus its dimension is \(r\) plus the number of even roots.
The even roots form a root system in their real span: they
are closed under reflections in their own roots, because
\(s_\alpha\beta=\beta-\beta(h_\alpha)\alpha\)
still has even parity. Its simple roots are exactly the
positive even roots that are not sums of two positive
even roots. The missing rank is the center dimension of
the fixed reductive algebra.

Here are explicit exhaustive procedures for both the orbits
and this root-system calculation. They use only integers.
Represent a root by its simple-root coefficient tuple \(c\).
Start with all unit tuples, close the set under
\[
 c_i\longmapsto c_i-\sum_j a_{ij}c_j
 \quad\text{with all other coordinates unchanged}.
 \tag{10.3}\label{eq:root-reflection-enumeration}
\]
This terminates at the known finite root system and obtains
every root: each root is in a Weyl orbit of a simple root,
as proved in Chapter~\ref{ch:root-systems}.
Independently, start with each as-yet-unseen binary tuple \(b\)
and close it under \eqref{eq:vogan-reflection}.
For each orbit retain its least integer representative
\(\sum_i2^{i-1}b_i\).
The resulting orbit sizes are respectively
\[
 (1,3),\quad(1,3,12),\quad(1,27,36),\quad
 (1,63,36,28),\quad(1,135,120).
\]
Their sums are \(2^2,2^4,2^6,2^7,2^8\).
For clarity, the core enumeration can be reproduced by
the following elementary pseudocode; \(\operatorname{closure}\)
adds images until no new tuple remains.
\[
\begin{aligned}
 \Phi&=\operatorname{closure}
       (\{e_1,\ldots,e_r\},\{s_1,\ldots,s_r\}),\\
 U&=\{0,1\}^r,\\
 \text{while }U\ne\varnothing:\quad
 b&=\operatorname{least}(U),\quad
 O=\operatorname{closure}(\{b\},\{s_1,\ldots,s_r\}),\\
 U&\leftarrow U\setminus O,\qquad
 \Phi_b=\{c\in\Phi:\sum b_i c_i\equiv0\pmod2\}.
\end{aligned}
\]
Here \(s_i\) in the first line is \eqref{eq:root-reflection-enumeration}
and in the orbit line is \eqref{eq:vogan-reflection}.
To check the fixed type without any Lie algebra identification
being assumed, take the positive elements of \(\Phi_b\),
delete every one that is a sum of two of them, and compute
their Cartan entries by
\[
 \frac{2c^T G d}{c^T G c},\qquad
 G=\operatorname{diag}(d_i)A,\quad
 (d_i)=
 \begin{cases}
 (1,3)&G_2,\\
 (1,1,2,2)&F_4,\\
 (1,\ldots,1)&E_r.
 \end{cases}
\]
This gives, in the table order, the root diagrams
\[
\begin{gathered}
 G_2,\ 2A_1;\qquad F_4,\ B_4,\ C_3+A_1;\\
 E_6,\ D_5,\ A_5+A_1;\qquad
 E_7,\ D_6+A_1,\ A_7,\ E_6;\qquad
 E_8,\ D_8,\ E_7+A_1 .
\end{gathered}
\]
The ranks in the two \(D_5\) and \(E_6\) rows are one less
than their ambient ranks, accounting for the central
\(\mathbb C\) in those rows. Counting the roots using the
same procedure gives ambient root counts
\(12,48,72,126,240\); the fixed dimensions then are exactly
those displayed. This procedure specifies every loop,
initial set, arithmetic operation, and stopping test in
the finite calculation, so the orbit and fixed-type
columns can be checked without using the real-form list
as input.
Each reflection is an actual conjugacy by
Proposition~\ref{res:J06}, so the calculation bounds the
number of inner forms from above. Different rows of one
type have different fixed dimensions, hence different
Killing signatures by Theorem~\ref{res:J04}; they cannot
be conjugate. Thus the bound is attained.

Only \(E_6\) has a nontrivial exceptional diagram involution.
It exchanges \(1\leftrightarrow5\), \(2\leftrightarrow4\)
and fixes \(3,6\). Reflection at a black fixed vertex
flips the other fixed vertex. Hence the four possible
colorings have the two classes \(00\) and
\(\{10,01,11\}\).
To identify the first, put
\[
 \beta_1=\tfrac12(\alpha_1+\alpha_5),\quad
 \beta_2=\tfrac12(\alpha_2+\alpha_4),\quad
 \beta_3=\alpha_3,\quad \beta_4=\alpha_6.
\]
For each paired orbit take the sums of its two \(e\)'s,
two \(f\)'s and two \(h\)'s; for each fixed vertex take
the individual generators. Each resulting
\((E_i,F_i,H_i)\) is an \(\mathfrak{sl}_2\) triple:
the roots in each paired orbit are nonadjacent, so all
cross brackets vanish. Direct evaluation on the simple
generators gives the Cartan matrix \(F_4\) in this order,
and \([F_i,E_j]=0\) for \(i\ne j\).
In the finite-dimensional adjoint module of the \(i\)-th
triple, \(E_j\) is therefore a lowest vector of weight
\(a_{ij}\). The rank-one classification gives
\((\operatorname{ad}E_i)^{1-a_{ij}}E_j=0\);
the negative relations follow in the same way.
All Serre relations have now been checked.
Theorem~\ref{res:F03} gives a nonzero homomorphism from
the simple \(F_4\) Serre algebra into the fixed algebra.
Its kernel is an ideal and cannot be the whole algebra,
so the homomorphism is injective and supplies a
\(52\)-dimensional fixed subalgebra. Direct application of
\eqref{eq:root-reflection-enumeration} shows that \(24\)
of the \(72\) roots of \(E_6\) are fixed by the diagram
involution. The other \(48\) roots form \(24\) pairs,
each contributing one fixed root vector. Thus the
fixed algebra has dimension at most \(4+24+24=52\).
Equality proves that it is precisely \(F_4\), and also
that the involution is \(+1\) on each fixed root line.

For the second class choose vertex \(3\) black. A fixed
root has an even coefficient at vertex \(3\) in exactly
eight of these \(24\) cases, by
\eqref{eq:root-reflection-enumeration}. The paired roots
still contribute \(24\) fixed vectors, so the fixed
algebra has dimension \(4+8+24=36\).
Its roots relative to the fixed Cartan are the projections
of the paired roots and these eight retained fixed roots.
The positive indecomposable roots in the \(\beta\) coordinates are
\[
 \beta_2,\quad\beta_1,\quad\beta_2+\beta_3,\quad\beta_4.
\]
Their squared lengths are \(1,1,1,2\), and their
successive inner products are
\(-\tfrac12,-\tfrac12,-1\), with all other inner
products zero. Their Cartan matrix is \(C_4\).
This verifies the fixed type as well as its dimension.
The two signatures are \(-26\) and \(6\), so these
two outer rows are distinct and are also distinct
from all inner rows.

Finally the split real algebra has signature equal to
the complex rank. Indeed its real Chevalley Cartan
has positive Killing form, and each real plane
\(\mathbb R e_\alpha+\mathbb R e_{-\alpha}\) contributes
one positive and one negative square. This identifies
the rows of signatures \(2,4,6,7,8\) as the split
forms. All remaining rows are constructed by their
explicit sign involutions. Combined with
Proposition~\ref{res:J05}, this proves existence,
exhaustiveness, and nonisomorphism of the exceptional
list. Its entries agree with
\cite[Section 42]{etingof} and
\cite[Theorem 6.105]{knapp}.
\end{proof}

\section{What the hypotheses prevent}

\begin{proposition}[Representative field and characteristic boundaries]
\label{res:J09}
The characteristic-zero, splitting, and finite-dimensionality
conditions in the preceding structure and representation
theorems cannot in general be removed.
\end{proposition}
\begin{proof}
\emph{A nonsplit real action.}
The one-dimensional abelian real algebra acting on
\(\mathbb R^2\) by
\(\left(\begin{smallmatrix}0&-1\\1&0\end{smallmatrix}\right)\)
has no invariant real line. After complexification it
has two eigenlines. Thus Lie's simultaneous-triangularization
conclusion needs a splitting hypothesis on the field.
Likewise a Cartan in \(\mathfrak{su}(2)\) has imaginary
nonzero root eigenvalues; its real algebra has the same
complex Dynkin type \(A_1\) as
\(\mathfrak{sl}_2(\mathbb R)\), but its Killing form is
negative definite rather than split.

\emph{An infinite-dimensional module for a semisimple algebra.}
In the Verma module \(M(0)\) of \(\mathfrak{sl}_2\), the
vectors \(f^jv\), \(j\geq0\), are independent by PBW.
The quotient sending \(v\) to \(1\) and all \(f^jv\),
\(j>0\), to zero is the trivial simple module.
It has no invariant section: a vector of weight zero
must be a multiple of \(v\), and \(fv\ne0\).
Hence this module is not completely reducible.
This does not contradict Weyl's theorem, which assumes
finite dimension.

\emph{Lie's theorem in positive characteristic.}
Let \(k\) have characteristic \(p>0\), and act on
\(V=k[t]/(t^p)\) by \(D=d/dt\), \(T=t\), and \(I\).
One has \([D,T]=I\), including on \(t^{p-1}\), and
\([D,I]=[T,I]=0\). Thus the three-dimensional
Heisenberg algebra is represented on \(V\).
If a nonzero vector were a common eigenvector for
\(D,T,I\), their commutator would kill it, whereas
\(I\) fixes it. There is no common eigenvector,
despite the acting algebra being nilpotent and
therefore solvable. Notice that \(I\) is not
nilpotent; Engel's theorem's hypothesis has not
been satisfied.

\emph{Cartan's Killing criterion in positive characteristic.}
For \(p>3\) let
\(W=\operatorname{Der}(k[t]/(t^p))\), with basis
\(e_i=t^{i+1}D\), \(-1\leq i\leq p-2\).
Its bracket is
\[
 [e_i,e_j]=(j-i)e_{i+j},
\]
with the term zero outside the indicated basis range.
This algebra is simple. Any nonzero ideal is stable
under \(\operatorname{ad}e_0\), whose eigenvalues
\(-1,0,\ldots,p-2\) are distinct in \(k\);
polynomial projections therefore put some \(e_i\)
in the ideal. Repeated bracketing with \(e_{-1}\)
puts \(e_{-1}\) in it, since the coefficients
\((i+1),i,\ldots,1\) are nonzero. Bracketing this
with the whole algebra gives
\(e_{-1},e_0,\ldots,e_{p-3}\). Finally
\([e_1,e_{p-3}]=(p-4)e_{p-2}\) supplies the last
vector. This proves simplicity.

Its Killing form is nevertheless zero.
Invariance under \(e_0\) says that \(B(e_i,e_j)\)
vanishes unless \(i+j=0\) modulo \(p\).
If the integer sum is \(p\), the product
\(\operatorname{ad}e_i\operatorname{ad}e_j\)
has zero diagonal: its final index has increased
by \(p\), beyond the basis range. The remaining
possibilities have integer sum zero. Here
\[
 B(e_0,e_0)=\sum_{a\in\mathbb F_p}a^2=0
 \qquad(p>3),
\]
and invariance with \([e_{-1},e_1]=2e_0\) gives
\(B(e_{-1},e_1)=-2B(e_0,e_0)=0\).
These exhaust the possibilities. Thus this simple
nonabelian algebra has zero Killing form.

For further positive-characteristic theory the appropriate
additional datum is a restricted operation \(x\mapsto x^{[p]}\):
it satisfies
\(\operatorname{ad}(x^{[p]})=(\operatorname{ad}x)^p\),
\((a x)^{[p]}=a^p x^{[p]}\), and the additive correction
given by the \(p\)-th-power commutator expansion.
For a derivation algebra the operation is the \(p\)-fold
composition, which is again a derivation by the binomial
formula in characteristic \(p\).
The restricted enveloping algebra is the quotient by
\(x^p-x^{[p]}\); these are additional relations, not the
ordinary characteristic-zero enveloping algebra.
The classification of restricted or modular simple
algebras, their modules, and the small-characteristic
exceptions belongs to the separate positive-characteristic
module in the scope map. The concrete counterexamples
above establish the boundaries used in this volume;
compare \cite[Section 15.4]{etingof},
\cite[Section 1]{premet}, and \cite[Sections A.1--A.2]{jantzen}.
\end{proof}
